\section{Введение}
Тематические модели описывают документы смесями тем, а темы --- распределениями по словам~\cite{blei2003}. При обучении на небольшой коллекции эти распределения оцениваются по ограниченному числу вхождений и совместных встречаемостей слов. В работе рассматривается способ дополнить такие наблюдения семантической информацией предобученной языковой модели. Базовой моделью служит AARTM, которая сочетает аддитивную регуляризацию~\cite{vorontsov2015} с локальным вниманием к соседним словам. Наша цель --- улучшить её предсказания без дообучения внешнего кодировщика и изменения набора обучаемых параметров AARTM.

Внешняя семантика уже используется в нескольких семействах тематических моделей. ETM задаёт вероятности слов через эмбеддинги слов и тем~\cite{dieng2020}. CombinedTM включает контекстное представление документа во вход вариационной сети~\cite{bianchi2021}. В FASTopic фиксированные представления документов связываются с обучаемыми представлениями тем и слов посредством оптимального транспорта~\cite{wu2024}. Эти подходы различаются тем, на каком этапе и в какой форме семантические представления участвуют в обучении.

Контекст можно учитывать и на уровне отдельных вхождений слов. CWTM использует контекстные представления слов в нейронной тематической модели~\cite{fang2024}; Contextual-Top2Vec на их основе выделяет темы и соответствующие им участки документов~\cite{angelov2024}. Мы используем представления фрагментов исходного текста и агрегируем их тематические назначения по словам словаря. Так получается матрица $Q=(q(t\mid w))$, которую можно передать AARTM как фиксированное семантическое ограничение. Агрегация сохраняет сведения о контекстах употребления, но усредняет различные значения одного слова.

По способу переноса знаний работа близка к дистилляции нейронных тематических моделей, предложенной Hoyle и соавторами~\cite{hoyle2020}. В их методе BERT-учитель задаёт распределения по словам для восстановления документов. Здесь учитель задаёт распределения по темам для каждого слова. Различается и способ его применения: одна матрица $Q$ используется сначала для инициализации AARTM, затем для KL-регуляризации её параметров. Вклад работы состоит в этом построении и его экспериментальной проверке на AARTM.

Предлагаемый метод, AARTM+Q, строит учителя только по обучающим документам. После построения $Q$ кодировщик больше не участвует в обновлениях тематической модели. Дополнительные затраты приходятся на кодирование фрагментов, кластеризацию и хранение матрицы размера $|V|\times T$. KL-согласование также может уменьшать разреженность тематических распределений. Поэтому наряду с полным методом мы рассматриваем более простой вариант, использующий только семантическую инициализацию.

Эксперимент должен установить, улучшает ли учитель восстановление скрытых слов и какая часть эффекта связана с инициализацией, а какая --- с KL. На пяти корпусах используются 10, 30 и 50 обучающих документов и три начальных состояния генератора случайных чисел. Методы сравниваются попарно на одних и тех же данных; качество предсказания оценивается по перплексии, согласованность наиболее вероятных слов тем --- по PMI и внешней NPMI. Среднее снижение document-completion PPL составляет 7,40\%, хотя внешняя NPMI систематически не растёт. Отдельная серия на 99\,049 документах arXiv исследует зависимость обучающей PMI от числа тем и сопоставляет AARTM+Q с FASTopic и CombinedTM.
